% ===== 原文 pp.26–27 =====

% [原文 p.26]
\begin{figure}[htbp]
\centering
\includegraphics[width=0.7\textwidth]{fig08.png}
\figtitle{8}{一个亏格 2、带三条有尖点的边界的曲面。}
\end{figure}

% [原文 p.26] 以下是第 ② 条 bullet 的续段（原文左缘右移约 27 pt），
% 故用 itemize + \item[] 恢复列表缩进；环境在本文件内闭合。
\begin{itemize}[leftmargin=1.6em,itemsep=2pt,topsep=2pt]
\item[]
有尖点的测地边界的曲面的标准模型（standard model），由粘接近乎紧的裤子得到——在每条边界上粘接一个近乎紧的裤子，使其边界长度与所粘边界的长度相匹配，并使其赤道点与该边界上的点对齐。该构造不涉及任何选择，因此关于标号同胚与 $\MCG$ 作用是自然的。

% [Tutor Note: p.26 的"标准模型曲面"（standard model surface）$\what{R}$ 是后续全部构造的几何载体：在每条测地线 $\beta_j$ 上粘接长度同为 $L_j$ 的近乎紧裤子，并让裤子的赤道点对齐该边界上的尖点（图 9），这个提升不含任何自由度。一个可核对的维数锚点：$\dim_\R\what{\cT}_{g,n}=\dim_\R\cT_g(L)+2n=(6g-6+2n)+2n=6g-6+4n$，它恰等于 $\dim_\R\cT_{g,2n}=6g-6+4n$，与原文 p.25 所说"The R dimension of $\what{\cT}_{g,n}$ is $2n$ greater than the R dimension of $\cT_g(L)$"一致（多出的 $2n$ 维正是 $n$ 个边界长度与 $n$ 个边界点位置）。这也解释了为什么开集 $\{L<c\}\subset\what{\cT}_{g,n}$ 能与 $\cT_{g,2n}/\prod_j\mathrm{Dehn}(\beta_j)$ 中的开集 $\{L<c\}$ 相互等价：Dehn 扭转是离散作用，不改变维数。]

\begin{figure}[htbp]
\centering
\includegraphics[width=0.5\textwidth]{fig09.png}
\figtitle{9}{对齐边界点与赤道点，以便粘合近乎紧的裤子。}
\end{figure}

对一个带标号、有尖点边界的曲面 $R$，记 $\what{R}$ 为粘接了近乎紧的裤子的标准模型曲面。$\what{R}$ 的穿刺点成对地带标号。由一般的双曲领口结果，长度小的测地线必然互不相交 [Bus92, Wlp10]。当长度 $L_1,\dots,L_n$ 适当小时，带标号的测地线 $\beta_1,\dots,\beta_n$ 在曲面 $\what{R}$ 上由"长度小且围出带标号的穿刺点"唯一确定。有尖点的边界、带标号曲面 $R$ \emph{等价}于带标号曲面 $\what{R}$ 模去关于 $\beta_j$ 的 Dehn 扭转（之所以是 Dehn 扭转，是因为边界点给在一个圆周上）。特别地，对适当小的 $c$，这一等价发生在 $\what{\cT}_{g,n}$ 的开子集 $\{L<c\}$ 与 $\cT_{g,2n}/\prod_j\mathrm{Dehn}(\beta_j)$ 中的开子集 $\{L<c\}$ 之间。

下一步是 $S^1$ 主丛的定义与等价性。讨论从定理 9 的短正合列开始，
\[
1\longrightarrow\mathrm{Dehn}(\beta)=\prod_j\mathrm{Dehn}(\beta_j)\longrightarrow\Stab(\beta)=\bigcap_j\Stab(\beta_j)\longrightarrow\MCG(R)\longrightarrow 1,
\]
% [Tutor Note: 与 p.21 定理 9 证明中的同型正合列对照可以读出这里的略写：那里第三条写的是 $\prod_{R(\gamma)\,\text{components}}\MCG(R')$——被切开曲面各分支的映射类群之积，而此处简记为 $\MCG(R)$。因此这条正合列的精确含义是 $1\to\prod_j\mathrm{Dehn}(\beta_j)\to\bigcap_j\Stab(\beta_j)\to\prod_{R(\beta)\ \text{分支}}\MCG\to 1$，与 p.27 首句"$\MCG(R)\approx\Stab(\beta)/\mathrm{Dehn}(\beta)$"一致。（原文 p.26 把它写成 $\MCG(R)$，属简记；译文照原样保留。）]
% [原文 p.27] 下面的"（此时…）"紧接上面的正合列，原文无段首缩进，故此处不空行。
（此时 $\Stab(\beta_j)=\Stab_0(\beta_j)$，因为保向的裤子同胚保持边界定向）。测地线 $\beta_1,\dots,\beta_n$ 定义了 $\what{\cT}_{g,n}$ 之上与 $\cT_{g,2n}/\mathrm{Dehn}(\beta)$ 之上的 Riemann 曲面丛（即万有曲线）的子集。只要 Riemann 曲面的自同构至多在单条测地线上作为旋转作用，这些子集就定义了定向圆周丛。小的长度 $L_1,\dots,L_n$ 与带标号的边界提供了这一条件。这些测地线定义了 Teichmüller 基上的圆周丛。我们在下面会看到，沿测地线的旋转定义了 $S^1$ 主结构。其次，由上述短正合列与标号的定义——测地线 $\beta\subset R$ 与 $\beta\subset\what{R}$ 之间的等价，以及圆周丛到基的投影——都与 $\MCG(R)\approx\Stab(\beta)/\mathrm{Dehn}(\beta)$ 的作用交换。这些测地线在商 $\{L<c\}\subset\what{\cT}_{g,n}/\MCG(R)\times(S^1)^n$（$\MCG(R)$ 与 $(S^1)^n$ 在 $\what{\cT}_{g,n}$ 上的作用彼此交换）与 $\{L<c\}\subset\cT_{g,2n}/\Stab(\beta)$ 上定义了等价的 orbifold $S^1$ 主丛。

$\cT_{g,2n}$ 的辛形式在开子集 $\{L<c\}$ 上定义了一个辛形式 $\omega_{\what{\cT}_{g,n}}$。这里用到 Fenchel--Nielsen 坐标与 $\dd\ell\wedge\dd\tau$ 公式。对任何包含多重曲线（multicurve）$\beta$ 的 $\what{R}$ 的裤子分解，形式 $\omega_{\what{\cT}_{g,n}}$ 都等于 $\sum_k\dd\ell_k\wedge\dd\tau_k$。重要的是，对带尖点边界 $\beta_j$ 的曲面 $R$，形式 $\omega_{\what{\cT}_{g,n}}$ 是由 $\dd\ell\wedge\dd\tau$ 公式的扩充解释给出的：每条边界各有一项，而现在的解释是 $\tau(\beta_j)$ 参数化所指定点的位置。参见图 3：参数 $\tau(\beta_j)$ 增大，对应于该点沿边界移动、且曲面内部位于右侧。辛形式的 $\MCG(R)$ 不变性是直接的。该形式到 $L$ 水平集上的限制、以及它在旋转边界点下的不变性，将在下面讨论。
\end{itemize}

% [Tutor Note: p.27 用"$\dd\ell\wedge\dd\tau$ 公式的扩充解释"来定义 $\omega_{\what{\cT}_{g,n}}$：对任何包含多重曲线 $\beta$ 的裤子分解都有 $\omega_{\what{\cT}_{g,n}}=\sum_k\dd\ell_k\wedge\dd\tau_k$，而每条边界 $\beta_j$ 多出一项、$\tau(\beta_j)$ 被解释为指定点在边界上的位置坐标。这正是命题 11 证明的要害：$L_j$ 固定后对应的 $\dd L_j$ 项消失，剩下的和就是"不含边界参数"的 $\omega_{\cT_g(L)}$ 公式；可与 p.25 对 $\cT_g(L)$ 的 $\sum_j\dd\ell_j\wedge\dd\tau_j$（"without boundary parameters"）逐项对照。]

\begin{itemize}[leftmargin=1.6em,itemsep=2pt,topsep=2pt]
\item $\cT_{g,n}$——带 $n$ 个穿刺点的带标号亏格 $g$ 双曲曲面构成的空间。$\cT_{g,n}$ 带有 $\MCG$ 不变的辛形式 $\omega$。$\cT_{g,n}$ 与 Teichmüller 空间 $\cT_g(0)$ 重合，后者中曲面的边界长度为零。
\end{itemize}

我们将把这三个辛流形联系起来。

% [Tutor Note: p.27 的三条 bullet 分属两页（①②在 p.25、③在此），原文用同一个 itemize；因 tr-03 的块在 p.25 末须闭合列表，此处另起一个 itemize 只装第 ③ 条。三条 bullet 依次是 $\cT_g(L)$、$\what{\cT}_{g,n}$、$\cT_{g,n}$ 三个辛流形，正是紧接着"这三个辛流形"所指的对象，读时可连起来看。]

\para{对 $\what{\cT}_{g,n}$ 的辛约化。}
我们考虑 FN 扭转、测地长度、尤其是矩映射的 Hamilton 几何
\[
\what{\cT}_{g,n}\overset{\mu}{\longrightarrow}\what{L}=(L_1^2/2,\dots,L_n^2/2)\in\R^n_{\ge 0}.
\]

记 $t_j$ 为边界 $\beta_j$ 上该点的单位速度无穷小旋转；$t_j$ 是 $\what{\cT}_{g,n}$ 上的向量场。就标准模型曲面而言，$t_j$ 就是 $\beta_j$ 的 FN 无穷小扭转向量场，而 $t_j$ 是第 $j$ 个近乎紧裤子的无穷小旋转。由扭转--长度对偶，有 $\omega_{\what{\cT}_{g,n}}(-t_j,\ )=\dd L_j$，且按比例缩放后的 $-L_jt_j$ 是单位无穷小旋转（单位时间流即一整圈旋转）。函数 $\half L_j^2$ 是相应的 Hamilton 势，因为
$\omega_{\what{\cT}_{g,n}}(-L_jt_j,\ )=$
